\chapter{Advantages and Limitations}
\setcounter{equation}{0}
\section{Advantages}
\begin{itemize}
\item \textbf{Efficient training:} XGBoost trained in 0.95 seconds in the reported benchmark, compared with 12.01 seconds for the neural network.
\item \textbf{Structured data suitability:} The method works with engineered URL and domain features and does not require a deep-learning stack or GPU for the described deployment.
\item \textbf{Feature interpretation:} Feature-importance scores help examine which URL characteristics influence the model.
\item \textbf{Missing-value handling:} XGBoost can handle missing values in extracted features.
\item \textbf{Browser use case:} The proposed pipeline is designed for URL assessment and risk display within a browser plug-in.
\item \textbf{Evaluation breadth:} Cross-dataset testing exposes limitations that a single dataset evaluation may hide.
\end{itemize}

\section{Limitations}
\begin{itemize}
\item \textbf{Dataset shift:} F1 falls from 91.16\% on the custom test set to 73.25\% on GramBeddings and 69.37\% on PhiUSIIL.
\item \textbf{Memory and inference trade-off:} XGBoost used more memory and had slower inference than the neural network in the reported comparison, despite faster training.
\item \textbf{Feature dependence:} The importance of HTTPS and other signals differs across datasets.
\item \textbf{Deployment scope:} The reported demonstration is limited to a single service and Chrome browser.
\item \textbf{Update requirements:} The described system does not include an automatic retraining pipeline to adapt continuously to changing phishing tactics.
\item \textbf{Residual risk:} No classifier can guarantee detection of every malicious URL; false positives and false negatives remain possible.
\end{itemize}

The paper's results support practical potential but also show that deployment should include monitoring, maintenance, and independent evaluation.
